\section{Introduction}\label{sec:introduction}

The scaling limit of a random surface does not by itself determine the
scaling limit of a walk on that surface. A time-parametrized walk also
depends on the electrical energy and on the mass assigned to vertices.
We study this additional limit for critical Fortuin--Kasteleyn (FK)--Ising
maps, using the common-coordinate and electrical convergence results of
the spectral companion~\cite{SC}.

\subsection{The model and the result}

Let $\mathcal F_n$ be the finite set of rooted pairs $(M,A)$, where $M$
is a connected planar multigraph with $n$ edges embedded in the oriented
sphere and $A\subseteq E(M)$. Rooting is at an oriented edge, and pairs
are identified by orientation-preserving homeomorphisms preserving the
root and $A$. Loops and multiple edges are allowed. If $M^*$ is the dual
map, set
\[
 A^*=\{e^*:e\notin A\},\qquad
 \ell(M,A)=k(A)+k(A^*)-1,
\]
where each component count includes isolated vertices of the corresponding
spanning subgraph. The critical spherical FK--Ising law is
\begin{equation}\label{eq:fk-law}
 \PP\bigl((M_n,A_n)=(M,A)\bigr)
 =\frac{2^{\ell(M,A)/2}}{Z_n},\qquad
 Z_n=\sum_{(M,A)\in\mathcal F_n}2^{\ell(M,A)/2}.
\end{equation}
Write $V_n=V(M_n)$, let $d_n$ be graph distance with every original edge
of length one, and let
\begin{equation}\label{eq:corner-measure}
 \mu_n(v)=\frac{\deg_{M_n}(v)}{2n}.
\end{equation}
A loop contributes two to the degree. Conditional on $(M_n,A_n)$, a
simple random walk chooses an incident half-edge uniformly and traverses
its edge. Its initial vertex has law $\mu_n$, which is stationary.
The walk uses all map edges, not just those in $A_n$.

Let $(S,h)$ be the ordinary unit-area $\sqrt3$-Liouville quantum gravity
(LQG) sphere: the quantum-sphere law disintegrated at area one, with its
distinguished points forgotten. Write $\mu_h$ for its area probability
measure and $D_h$ for its intrinsic metric in a fixed deterministic
normalization. Liouville Brownian motion is the diffusion associated with
the closure on $L^2(\mu_h)$ of
\begin{equation}\label{eq:continuum-form}
 \E_h(f,f)=\frac12\int_S \abs{\nabla_g f}^2\dd\mathrm{vol}_g,
 \qquad f\in C^\infty(S),
\end{equation}
where $g$ is any smooth metric in the conformal class of the sphere.
Conformal invariance of the two-dimensional Dirichlet integral makes
this definition independent of $g$. The factor $1/2$ fixes the time
convention.

For a walk $X^n$, let $\overline X^n$ be its continuous interpolation on
the cable realization $\overline M_n$, in which each edge is a unit
interval and each selected edge is traversed linearly in one unit of
time. We also write $d_n$ for the cable path metric. We use the
curve-decorated Gromov--Hausdorff--Prokhorov topology: after isometric
embeddings into common metric spaces, the underlying spaces converge in
Hausdorff distance, the measures in Prokhorov distance, and every curve
uniformly on each bounded time interval.

\begin{theorem}\label{thm:main}
Let $(M_n,A_n)$ have law \eqref{eq:fk-law}, and let $a_n>0$ be any
deterministic sequence with $a_n\to0$ such that
\begin{equation}\label{eq:assumed-ghp}
 (V_n,a_nd_n,\mu_n)\ \Longrightarrow\ (S,D_h,\mu_h)
\end{equation}
in Gromov--Hausdorff--Prokhorov law. Given the map, start $r$ independent
simple random walks $X^{n,1},\ldots,X^{n,r}$ independently from $\mu_n$.
Given $h$, start $r$ independent Liouville Brownian motions
$B^{h,1},\ldots,B^{h,r}$ independently from $\mu_h$, with the convention
\eqref{eq:continuum-form}. For every fixed integer $r\ge1$,
\begin{equation}\label{eq:main-limit}
 \left(\overline M_n,a_nd_n,\mu_n;
       (\overline X^{n,j}_{nt})_{t\ge0,\,1\le j\le r}\right)
 \ \Longrightarrow\
 \left(S,D_h,\mu_h;
       (B^{h,j}_t)_{t\ge0,\,1\le j\le r}\right)
\end{equation}
in the curve-decorated Gromov--Hausdorff--Prokhorov topology.
\end{theorem}

Thus the deterministic acceleration can be taken to be exactly $n$.
The finite-product assertion retains the conditional walk law in the
random environment, rather than only the distribution of one trajectory
after averaging over the surface.

The proof uses more than the metric-measure premise
\eqref{eq:assumed-ghp}. The spectral companion supplies a joint conformal
coordinate, unit conductivity, uniform local energy recovery, and
convergence of centered inverse operators. We state these inputs
precisely in Proposition~\ref{prop:inputs}; their proofs and the
foundational companion results on which they depend are not reproduced
here. The path-convergence argument from those inputs is proved in full.

\subsection{History and related work}

Sheffield's inventory-accumulation encoding related critical FK-decorated
planar maps to correlated Brownian paths~\cite{Sheffield2016Inventory}.
Gwynne and Sun proved that the finite-volume FK contour encoding converges
to correlated Brownian motion conditioned to stay in the quadrant and
return to the origin~\cite[Theorem~1.8]{GwynneSun2015Finite}.
The mating-of-trees theory of Duplantier, Miller, and Sheffield, and its
finite-volume sphere form due to Miller and Sheffield, identify the
corresponding continuum quantum
surfaces~\cite{DuplantierMillerSheffield2021Mating,MillerSheffield2019SphereMating}.
The Brownian paths in that encoding describe boundary lengths; the vertex
walk considered here is a different stochastic process.
Liouville Brownian motion was constructed by Garban, Rhodes, and Vargas
and by Berestycki~\cite{GarbanRhodesVargas2016LBM,Berestycki2015Diffusion}.
Its Dirichlet-form description is developed in
\cite[Theorem~1.7]{GarbanRhodesVargas2014Dirichlet}.

For mated-CRT maps, Gwynne, Miller, and Sheffield proved quenched walk
convergence modulo time
parametrization~\cite[Theorem~3.4]{GwynneMillerSheffield2021Tutte}.
Berestycki and Gwynne subsequently retained the
parametrization~\cite[Theorems~1.2 and~1.4]{BerestyckiGwynne2022MatedCRTLBM},
using a deterministic median-exit clock comparable to inverse cell area.
Its asymptotic constant is not identified there
\cite[Remark~1.3]{BerestyckiGwynne2022MatedCRTLBM}.
The setting here is instead the ordinary finite FK--Ising sphere, the
walks start in stationarity, and the limit is joint with the intrinsic
metric and area measure. The exact clock in Theorem~\ref{thm:main} rests
on the unit-conductivity theorem of the spectral
companion~\cite[Theorem~5.2]{SC}.

\subsection{From operators to paths}

The principal issue is tightness with the time parameter retained.
Convergence of eigenvalues, or even convergence of the finite-dimensional
walk laws, does not exclude rapid excursions of appreciable spatial
diameter. The companion's convergence of centered inverses gives
positive-time semigroup convergence. Uniform tracking of the exploration
tour, which lists the graph corners and parametrizes quantum area in the
limit, then identifies the finite-dimensional distributions in one common
coordinate.

To control entire paths, we use smooth spatial cutoffs. Local energy
recovery provides uniformly accurate graph cutoffs with asymptotically
no excess energy. Resolvent smoothing gives test functions whose
generators are uniformly bounded at each fixed smoothing parameter.
A reversible forward--backward martingale decomposition converts the
form-norm error between these two approximations into an error uniform
along a stationary trajectory. Optional sampling of the smoothed
cutoffs then gives a stopping-time tightness estimate.

The local contribution is this passage from the companion's electrical
and operator limits to stationary, time-parametrized path convergence.
The maximal inequality is a finite-state analogue of the classical
Lyons--Zheng forward--backward
decomposition~\cite{LyonsZheng1988}; we prove the jump-chain identity
directly. The last tightness step uses Aldous's
criterion~\cite{Aldous1978}.
Neither an additional mixing assumption nor a heat-kernel estimate is
needed for this passage. Stationarity is essential to the argument as
given; we do not assert convergence from arbitrary starting vertices.

Section~\ref{sec:inputs} records the precise companion inputs and checks
the normalization. Section~\ref{sec:operators} obtains the
finite-dimensional laws. Section~\ref{sec:tightness} proves the stationary
maximal estimate and path tightness. Section~\ref{sec:paths} transfers the
limit to cable paths, removes the exponential clock, and assembles the
joint convergence in random environments.
